% ============================================================
% Section 141: 光照下半导体的准费米能级
% ============================================================
\section{半导体物理：光照下半导体的准费米能级}
\label{sec:quasi-fermi-levels}

\begin{modelbox}[题目：光激发n型硅的准费米能级]
一块含有施主杂质浓度 $N_D=1\times10^{17}\,\text{cm}^{-3}$ 的半导体硅，在光照射下体内均匀产生 $\Delta n=\Delta p=1\times10^{10}\,\text{cm}^{-3}$ 的光生载流子。试画出此时硅的能带图，并准确标出准费米能级的位置。
\end{modelbox}

\begin{tipbox}[考点快速分析]
\begin{itemize}
    \item \textbf{热平衡}：$n_0\approx N_D$，$p_0=n_i^2/n_0$
    \item \textbf{小注入}：$\Delta n\ll n_0$，多子浓度几乎不变，少子浓度剧增
    \item \textbf{准费米能级}：电子 $E_{Fn}$、空穴 $E_{Fp}$，分别由各自的载流子浓度定义
    \item \textbf{分裂程度}：$E_{Fn}-E_{Fp}$ 表征系统偏离热平衡的程度
\end{itemize}
\end{tipbox}

\begin{derivationbox}[标准解题过程]
\textbf{1. 热平衡时的载流子浓度}

已知硅室温 $n_i=1.5\times10^{10}\,\text{cm}^{-3}$。施主全电离下：
\[
n_0\approx N_D=1\times10^{17}\,\text{cm}^{-3},\qquad p_0=\frac{n_i^2}{n_0}=\frac{2.25\times10^{20}}{10^{17}}=2.25\times10^3\,\text{cm}^{-3}.
\]
热平衡费米能级 $E_F$ 靠近导带底。

\textbf{2. 光照下的非平衡载流子浓度}

光注入后：
\[
n=n_0+\Delta n\approx10^{17}+10^{10}\approx10^{17}\,\text{cm}^{-3}\approx n_0,
\]
\[
p=p_0+\Delta p\approx2.25\times10^3+10^{10}\approx10^{10}\,\text{cm}^{-3}.
\]
多子（电子）浓度几乎不变，少子（空穴）浓度剧增约 $4.4\times10^6$ 倍。

\textbf{3. 准费米能级的位置}

电子准费米能级：
\[
E_{Fn}=E_i+k_BT\ln\frac{n}{n_i}\approx E_i+0.026\times\ln(6.67\times10^6)\approx E_i+0.41\,\text{eV}.
\]
与热平衡 $E_F$ 基本重合，仍靠近导带底。

空穴准费米能级：
\[
E_{Fp}=E_i-k_BT\ln\frac{p}{n_i}\approx E_i-0.026\times\ln(0.67)\approx E_i+0.01\,\text{eV}.
\]
大幅下移至 $E_i$ 附近（禁带中央）。

\textbf{4. 能带图}

在能带图上：导带底 $E_c$、价带顶 $E_v$、本征费米能级 $E_i$ 为参考线。$E_{Fn}$ 靠近 $E_c$；$E_{Fp}$ 位于 $E_i$ 附近。两者显著分裂：
\begin{equation}
\boxed{E_{Fn}-E_{Fp}\approx0.4\,\text{eV}.}
\end{equation}
\end{derivationbox}

\begin{formulabox}[答案汇总]
\begin{center}
\begin{tabular}{ll}
\toprule
\textbf{物理量} & \textbf{位置} \\
\midrule
电子准费米能级 $E_{Fn}$ & 靠近导带底，与热平衡 $E_F$ 基本重合 \\
空穴准费米能级 $E_{Fp}$ & 大幅下移至禁带中央（$E_i$ 附近） \\
能带图特征 & $E_{Fn}$ 与 $E_{Fp}$ 显著分裂，间距约 $0.4\,\text{eV}$ \\
\bottomrule
\end{tabular}
\end{center}
\end{formulabox}

\begin{insightbox}[物理图像]
\begin{itemize}
    \item 小注入条件下（$\Delta n\ll n_0$），多子准费米能级几乎不变，少子准费米能级发生显著移动。
    \item 准费米能级的分裂程度 $E_{Fn}-E_{Fp}$ 定量表征了系统偏离热平衡的程度，是分析 pn 结、太阳电池、光电探测器等器件非平衡输运的核心概念。
    \item 本题中空穴准费米能级移至 $E_i$ 附近，说明空穴浓度已被提升至接近本征水平，硅从强 n 型转变为近本征的少子分布。
\end{itemize}
\end{insightbox}
